\section{Homework Week 12}

\begin{exercise}
    Write out a proof of Maschke's Theorem in the case of representations over
    \(\mathbb{C}\) along the following lines.
    
    Given a representation
    \[
        \rho : G \longrightarrow \operatorname{GL}(V),
    \]
    where \(V\) is a vector space over \(\mathbb{C}\), let
    \[
        (\cdot,\cdot)
    \]
    be any positive definite Hermitian form on \(V\). Define a new form
    \[
        (\cdot,\cdot)_1
    \]
    on \(V\) by
    \[
        (v,w)_1
        =
        \frac{1}{|G|}
        \sum_{g\in G} (gv,gw).
    \]
    Show that \((\cdot,\cdot)_1\) is a positive definite Hermitian form, preserved
    under the action of \(G\); that is,
    \[
        (v,w)_1 = (gv,gw)_1
    \]
    always.
    
    If \(W\) is a subrepresentation of \(V\), show that
    \[
        V = W \oplus W^\perp
    \]
    as representations.
    \end{exercise}
    
    \begin{proof}
        Throughout, write \(gv\) for \(\rho(g)(v)\).
        
        First we show that \((\cdot,\cdot)_1\) is Hermitian. Since each map
        \[
            v \longmapsto gv
        \]
        is linear, and since \((\cdot,\cdot)\) is Hermitian, the averaged form
        \[
            (v,w)_1
            =
            \frac{1}{|G|}
            \sum_{g\in G}(gv,gw)
        \]
        is again sesquilinear and Hermitian.
        
        Next we prove positive definiteness. If \(v\neq 0\), then \(gv\neq 0\) for every
        \(g\in G\), because \(g\) acts invertibly on \(V\). Hence
        \[
            (gv,gv)>0
        \]
        for all \(g\in G\). Therefore
        \[
            (v,v)_1
            =
            \frac{1}{|G|}
            \sum_{g\in G}(gv,gv)
            >0.
        \]
        Thus \((\cdot,\cdot)_1\) is positive definite.
        
        Now let \(h\in G\). Then
        \[
            (hv,hw)_1
            =
            \frac{1}{|G|}
            \sum_{g\in G}(ghv,ghw).
        \]
        As \(g\) runs through \(G\), so does \(gh\). Hence, after the change of variable
        \(t=gh\), we obtain
        \[
            (hv,hw)_1
            =
            \frac{1}{|G|}
            \sum_{t\in G}(tv,tw)
            =
            (v,w)_1.
        \]
        Therefore \((\cdot,\cdot)_1\) is preserved by the action of \(G\).
        
        Let \(W\subseteq V\) be a subrepresentation, and let \(W^\perp\) denote the
        orthogonal complement of \(W\) with respect to \((\cdot,\cdot)_1\). Since
        \((\cdot,\cdot)_1\) is positive definite and \(V\) is finite-dimensional, we have
        a vector space direct sum
        \[
            V = W \oplus W^\perp.
        \]
        It remains to show that \(W^\perp\) is \(G\)-stable.
        
        Take \(u\in W^\perp\), \(h\in G\), and \(w\in W\). Since \(W\) is a subrepresentation,
        we have
        \[
            h^{-1}w\in W.
        \]
        Using the \(G\)-invariance of \((\cdot,\cdot)_1\), we get
        \[
            (hu,w)_1
            =
            (hu,h(h^{-1}w))_1
            =
            (u,h^{-1}w)_1
            =
            0,
        \]
        because \(u\in W^\perp\) and \(h^{-1}w\in W\). Hence \(hu\in W^\perp\).
        
        Thus \(W^\perp\) is a subrepresentation of \(V\). Consequently,
        \[
            V = W \oplus W^\perp
        \]
        as representations. This proves Maschke's theorem over \(\mathbb{C}\).
        \end{proof}

    \begin{exercise}
    Prove the following theorem of Burnside: let \(G\) be a finite group, let \(k\)
    be an algebraically closed field in which \(|G|\) is invertible, and let
    \[
        \rho : G \longrightarrow \operatorname{GL}(V)
    \]
    be a representation over \(k\). By taking a basis of \(V\), write each
    endomorphism \(\rho(g)\) as a matrix. Let
    \[
        \dim V = n.
    \]
    Show that the representation is simple if and only if there exist \(n^2\)
    elements
    \[
        g_1,\ldots,g_{n^2}\in G
    \]
    so that the matrices
    \[
        \rho(g_1),\ldots,\rho(g_{n^2})
    \]
    are linearly independent, and that this happens if and only if the algebra
    homomorphism
    \[
        kG \longrightarrow \operatorname{End}_k(V)
    \]
    is surjective. Note that \(\rho\) itself is generally not surjective.
    \end{exercise}

    \begin{proof}
        Let
        \[
            \Phi : kG \longrightarrow \operatorname{End}_k(V)
        \]
        be the algebra homomorphism induced by the representation, namely
        \[
            \Phi\left(\sum_{g\in G} a_g g\right)
            =
            \sum_{g\in G} a_g \rho(g).
        \]
        Set
        \[
            A := \operatorname{Im}\Phi.
        \]
        Then \(A\) is the \(k\)-subalgebra of \(\operatorname{End}_k(V)\) spanned by the
        matrices
        \[
            \rho(g), \qquad g\in G.
        \]
        
        A subspace \(W\subseteq V\) is \(G\)-stable if and only if it is stable under every
        element of \(A\). Hence the representation \(V\) is simple as a representation of
        \(G\) if and only if \(V\) is a simple left \(A\)-module.
        
        Since \(|G|\) is invertible in \(k\), Maschke's theorem implies that \(kG\) is
        semisimple. Therefore its quotient algebra \(A\) is also semisimple. Since \(k\) is
        algebraically closed and \(A\) is finite-dimensional semisimple, the Artin--Wedderburn
        theorem gives
        \[
            A \cong
            M_{m_1}(k)\oplus \cdots \oplus M_{m_r}(k).
        \]
        
        Suppose first that \(V\) is simple as an \(A\)-module. Since \(A\subseteq
        \operatorname{End}_k(V)\), the action of \(A\) on \(V\) is faithful. But for the algebra
        \[
            M_{m_1}(k)\oplus \cdots \oplus M_{m_r}(k),
        \]
        a simple module is supported on exactly one simple matrix-algebra summand. If \(r>1\),
        then all other summands act by zero, contradicting faithfulness. Hence \(r=1\), and
        therefore
        \[
            A \cong M_m(k)
        \]
        for some \(m\).
        
        The unique simple left \(M_m(k)\)-module is \(k^m\), so
        \[
            \dim_k V = m.
        \]
        Since \(\dim_k V=n\), we have \(m=n\). Hence
        \[
            \dim_k A = m^2 = n^2.
        \]
        But \(A\subseteq \operatorname{End}_k(V)\), and
        \[
            \dim_k \operatorname{End}_k(V)=n^2.
        \]
        Therefore
        \[
            A=\operatorname{End}_k(V).
        \]
        Thus \(\Phi:kG\to \operatorname{End}_k(V)\) is surjective.
        
        Conversely, suppose that
        \[
            \Phi:kG\longrightarrow \operatorname{End}_k(V)
        \]
        is surjective. Then \(A=\operatorname{End}_k(V)\). Let \(0\neq W\subseteq V\) be an
        \(A\)-submodule. Choose \(0\neq w\in W\). For any \(v\in V\), there exists a linear
        endomorphism \(T\in \operatorname{End}_k(V)\) such that
        \[
            T(w)=v.
        \]
        Since \(A=\operatorname{End}_k(V)\), we have \(T\in A\). Because \(W\) is stable under
        \(A\), it follows that
        \[
            v=T(w)\in W.
        \]
        Thus \(W=V\). Hence \(V\) is simple.
        
        So the representation is simple if and only if
        \[
            \Phi:kG\longrightarrow \operatorname{End}_k(V)
        \]
        is surjective.
        
        It remains to relate this to the existence of \(n^2\) linearly independent matrices
        \(\rho(g_i)\). Since
        \[
            A=\operatorname{Span}_k\{\rho(g)\mid g\in G\},
        \]
        there exist elements
        \[
            g_1,\ldots,g_{n^2}\in G
        \]
        such that
        \[
            \rho(g_1),\ldots,\rho(g_{n^2})
        \]
        are linearly independent if and only if
        \[
            \dim_k A \geq n^2.
        \]
        But \(A\subseteq \operatorname{End}_k(V)\), and
        \[
            \dim_k \operatorname{End}_k(V)=n^2.
        \]
        Hence this happens if and only if
        \[
            \dim_k A=n^2,
        \]
        which is equivalent to
        \[
            A=\operatorname{End}_k(V).
        \]
        Therefore this happens if and only if \(\Phi\) is surjective.
        
        Combining the previous equivalences, the representation is simple if and only if there
        exist \(n^2\) elements
        \[
            g_1,\ldots,g_{n^2}\in G
        \]
        such that
        \[
            \rho(g_1),\ldots,\rho(g_{n^2})
        \]
        are linearly independent, and this is equivalent to the surjectivity of
        \[
            kG\longrightarrow \operatorname{End}_k(V).
        \]
        \end{proof}
    
    \begin{exercise}
    Let
    \[
        G = \langle x \mid x^n = 1 \rangle
    \]
    be a cyclic group of order \(n\), and let \(\zeta_n\in\mathbb{C}\) be a primitive
    \(n\)-th root of unity. Then the simple complex characters of \(G\) are the
    \(n\) functions
    \[
        \chi_r(x^s) = \zeta_n^{rs},
    \]
    where
    \[
        0 \leq r \leq n-1.
    \]
    \end{exercise}

    \begin{proof}
        Let
        \[
            G=\langle x\mid x^n=1\rangle.
        \]
        For each \(0\leq r\leq n-1\), define a one-dimensional representation
        \[
            \rho_r:G\longrightarrow \operatorname{GL}_1(\mathbb{C})
        \]
        by
        \[
            \rho_r(x)=\zeta_n^r.
        \]
        Since \((\zeta_n^r)^n=1\), this is a well-defined representation of \(G\). Its character
        is given by
        \[
            \chi_r(x^s)
            =
            \operatorname{tr}(\rho_r(x^s))
            =
            \operatorname{tr}\bigl((\zeta_n^r)^s\bigr)
            =
            \zeta_n^{rs}.
        \]
        Each \(\rho_r\) is one-dimensional, hence simple.
        
        It remains to prove that every simple complex representation of \(G\) is one of these.
        Let
        \[
            \rho:G\longrightarrow \operatorname{GL}(V)
        \]
        be a simple complex representation, and write
        \[
            T:=\rho(x).
        \]
        Since \(x^n=1\), we have
        \[
            T^n=I.
        \]
        Therefore the minimal polynomial of \(T\) divides
        \[
            X^n-1.
        \]
        Over \(\mathbb{C}\), the polynomial \(X^n-1\) splits into distinct linear factors:
        \[
            X^n-1=\prod_{r=0}^{n-1}(X-\zeta_n^r).
        \]
        Hence \(T\) is diagonalizable, and
        \[
            V=\bigoplus_{\lambda^n=1} V_\lambda,
        \]
        where
        \[
            V_\lambda=\{v\in V\mid T v=\lambda v\}.
        \]
        
        Each eigenspace \(V_\lambda\) is \(G\)-stable, because \(G\) is generated by \(x\), and
        \(x\) acts on \(V_\lambda\) by the scalar \(\lambda\). Since \(V\) is simple, exactly one
        eigenspace is nonzero. Thus
        \[
            V=V_\lambda
        \]
        for some \(n\)-th root of unity \(\lambda\).
        
        Moreover, if \(\dim_{\mathbb{C}}V>1\), then any one-dimensional subspace of \(V\) is
        stable under \(G\), because \(x\) acts on all of \(V\) by the scalar \(\lambda\). This would
        contradict the simplicity of \(V\). Hence
        \[
            \dim_{\mathbb{C}}V=1.
        \]
        Therefore \(V\) is a one-dimensional representation on which
        \[
            x
        \]
        acts by some \(n\)-th root of unity \(\lambda\). Since \(\zeta_n\) is primitive, there is a
        unique \(r\) with
        \[
            0\leq r\leq n-1
        \]
        such that
        \[
            \lambda=\zeta_n^r.
        \]
        Thus \(V\cong \rho_r\), and its character is
        \[
            \chi_r(x^s)=\zeta_n^{rs}.
        \]
        
        Finally, the characters \(\chi_0,\ldots,\chi_{n-1}\) are distinct, since
        \[
            \chi_r(x)=\zeta_n^r
        \]
        are distinct for \(0\leq r\leq n-1\). Therefore the simple complex characters of \(G\)
        are precisely
        \[
            \chi_r(x^s)=\zeta_n^{rs},
            \qquad 0\leq r\leq n-1.
        \]
        \end{proof}